\documentclass{beamer}
\usepackage[utf8]{inputenc}

\usetheme{Madrid}
\usecolortheme{default}
\usepackage{amsmath,amssymb,amsfonts,amsthm}
\usepackage{txfonts}
\usepackage{tkz-euclide}
\usepackage{listings}
\usepackage{adjustbox}
\usepackage{array}
\usepackage{tabularx}
\usepackage{gvv}
\usepackage{lmodern}
\usepackage{circuitikz}
\usepackage{tikz}
\usepackage{graphicx}

\setbeamertemplate{page number in head/foot}[totalframenumber]

\usepackage{tcolorbox}
\tcbuselibrary{minted,breakable,xparse,skins}



\definecolor{bg}{gray}{0.95}
\DeclareTCBListing{mintedbox}{O{}m!O{}}{%
	breakable=true,
	listing engine=minted,
	listing only,
	minted language=#2,
	minted style=default,
	minted options={%
		linenos,
		gobble=0,
		breaklines=true,
		breakafter=,,
		fontsize=\small,
		numbersep=8pt,
		#1},
	boxsep=0pt,
	left skip=0pt,
	right skip=0pt,
	left=25pt,
	right=0pt,
	top=3pt,
	bottom=3pt,
	arc=5pt,
	leftrule=0pt,
	rightrule=0pt,
	bottomrule=2pt,
	toprule=2pt,
	colback=bg,
	colframe=orange!70,
	enhanced,
	overlay={%
		\begin{tcbclipinterior}
			\fill[orange!20!white] (frame.south west) rectangle ([xshift=20pt]frame.north west);
	\end{tcbclipinterior}},
	#3,
}
\lstset{
	language=C,
	basicstyle=\ttfamily\small,
	keywordstyle=\color{blue},
	stringstyle=\color{orange},
	commentstyle=\color{green!60!black},
	numbers=left,
	numberstyle=\tiny\color{gray},
	breaklines=true,
	showstringspaces=false,
}
%------------------------------------------------------------
%This block of code defines the information to appear in the
%Title page
\title %optional
{12.268}
%\subtitle{A short story}

\author % (optional)
{RAVULA SHASHANK REDDY - EE25BTECH11047}

 \begin{document}
	
	
	\frame{\titlepage}
	\begin{frame}{Question}
    A $3 \times 3$ matrix has elements such that its trace is $11$ and its determinant is $36$. The eigenvalues of the matrix are all known to be positive integers. Find the largest eigenvalue of the matrix.
\end{frame}
\begin{frame}{Solution}
    Let the eigenvalues of the $3 \times 3$ matrix $\vec{A}$ be $\lambda_1, \lambda_2, \lambda_3$.  

We are given:
\begin{align}
\text{trace}(\vec{A}) = \lambda_1 + \lambda_2 + \lambda_3 = 11,\\ 
\det(\vec{A}) = \lambda_1 \lambda_2 \lambda_3 = 36\\
\lambda_1, \lambda_2, \lambda_3 > 0
\end{align}

Check the possible factor triplets of 36 and their sums:
\begin{align}
1 + 4 + 9 &= 14 \not=11 \\
2 + 2 + 9 &= 13 \not=11\\
2 + 3 + 6 &= 11 \\
3 + 3 + 4 &= 10 \not=11
\end{align}
\end{frame}
\begin{frame}{Solution}
Hence, the eigenvalues are:
\begin{align}
\myvec{\lambda_1 \\ \lambda_2 \\ \lambda_3} = \myvec{2 \\ 3 \\ 6}.
\end{align}

Thus, the largest eigenvalue is:
\[
\lambda_{\max} = 6.
\]
\end{frame}
\begin{frame}[fragile]
\frametitle{C Code}
    \begin{lstlisting}
        #include <stdio.h>

int main() {
    int trace = 11;
    int det = 36;
    int i, j, k;

    printf("The positive integer eigenvalues are:\n");

    for(i = 1; i <= trace; i++) {           // possible first eigenvalue
        for(j = 1; j <= trace - i; j++) {   // second eigenvalue
            k = trace - i - j;              // third eigenvalue from trace
            if(k > 0 && i * j * k == det) { // check product = determinant
                printf("Eigenvalues: %d, %d, %d\n", i, j, k);
                \end{lstlisting}
                \end{frame}
\begin{frame}[fragile]
\frametitle{C Code}
    \begin{lstlisting}
                int max = i;
                if(j > max) max = j;
                if(k > max) max = k;
                printf("Largest eigenvalue: %d\n", max);
            }
        }
    }

    return 0;
}
    \end{lstlisting}
\end{frame}
\begin{frame}[fragile]
\frametitle{Python Code}
    \begin{lstlisting}
import numpy as np

trace = 11       # sum of eigenvalues
det = 36         # product of eigenvalues

valid_eigenvalues = []

# Loop over all possible positive integers for first two eigenvalues
for i in range(1, trace + 1):
    for j in range(1, trace - i + 1):
        k = trace - i - j
        if k > 0 and i * j * k == det:
            eigenvalues = np.array([i, j, k])
            valid_eigenvalues.append(eigenvalues)
            print("Eigenvalues:", eigenvalues)
            print("Largest eigenvalue:", np.max(eigenvalues))
if valid_eigenvalues:
    largest_overall = max(np.max(triplet) for triplet in valid_eigenvalues)
    print("Largest eigenvalue among all triplets:", largest_overall)        
    \end{lstlisting}
\end{frame}
\begin{frame}[fragile]
\frametitle{Python Shared}
    \begin{lstlisting}
        import ctypes

# Load the shared library
lib = ctypes.CDLL('./eigen.so')

# Define the function argument and return types
lib.find_largest_eigen.argtypes = [ctypes.c_int, ctypes.c_int]
lib.find_largest_eigen.restype = ctypes.c_int

# Inputs
trace = 11
det = 36

# Call the C function
largest = lib.find_largest_eigen(trace, det)

print(f"Largest eigenvalue for trace={trace}, det={det} is: {largest}")
    \end{lstlisting}
\end{frame}
\end{document}
